\documentclass{article}
\usepackage{mathtools} 
\usepackage{fontspec}
\usepackage[UTF8]{ctex}
\usepackage{amsthm}
\usepackage{mdframed}
\usepackage{xcolor}
\usepackage{amssymb}
\usepackage{amsmath}


% 定义新的带灰色背景的说明环境 zremark
\newmdtheoremenv[
  backgroundcolor=gray!10,
  % 边框与背景一致，边框线会消失
  linecolor=gray!10
]{zremark}{说明}

% 通用矩阵命令: \flexmatrix{矩阵名}{元素符号}{行数}{列数}
\newcommand{\flexmatrix}[4]{
  \[
  #1 = \begin{pmatrix}
    #2_{11}     & #2_{12}     & \cdots & #2_{1#4}   \\
    #2_{21}     & #2_{22}     & \cdots & #2_{2#4}   \\
    \vdots      & \vdots      & \ddots & \vdots     \\
    #2_{#31}    & #2_{#32}    & \cdots & #2_{#3#4}
  \end{pmatrix}
  \]
}

% 简化版命令（默认矩阵名为A，元素符号为a）: \quickmatrix{行数}{列数}
\newcommand{\quickmatrix}[2]{\flexmatrix{A}{a}{#1}{#2}}

\begin{document}
\title{注释 16.2}
\author{张志聪}
\maketitle

% \begin{zremark}
%   证明：
%   设 $f:\mathbb R^2\setminus\{(0,0)\}\to\mathbb R$。
%   \begin{align*}
%     \lim_{(x,y)\to(0,0)} f(x,y) = L
%   \end{align*}
%   当且仅当
%   \begin{align*}
%     L = \lim_{r\to0^+} f(r\cos\theta,r\sin\theta)
%   \end{align*}
%   其中
%   \[
%     x=r\cos\theta,\qquad y=r\sin\theta,\qquad r>0,\ \theta\in[0,2\pi).
%   \]
% \end{zremark}



\end{document}